\documentclass[11pt]{article}

\usepackage[margin=1in]{geometry}
\usepackage{amsmath,amssymb,amsthm,mathtools}
\usepackage{microtype}
\usepackage[hidelinks]{hyperref}

\newtheorem{theorem}{Theorem}[section]
\newtheorem{lemma}[theorem]{Lemma}
\newtheorem{corollary}[theorem]{Corollary}
\theoremstyle{definition}
\newtheorem{definition}[theorem]{Definition}

\newcommand{\Pp}{\mathbb{P}}
\newcommand{\Cov}{\operatorname{Cov}}
\newcommand{\1}{\mathbf{1}}
\newcommand{\Exp}{\operatorname{Exp}}

\hypersetup{pdftitle={Signed Faithfulness and Independence Collapse for Luce Choice Events}}

\begin{document}

\begin{center}
{\LARGE\bfseries Signed Faithfulness and Independence Collapse\\[3pt]
for Luce Choice Events\par}
\end{center}

\begin{abstract}
Let a Plackett--Luce ranking be generated by independent exponential clocks, and consider events
asserting that a specified item ranks first in a specified menu. We derive the exact covariance of
every pair of such events. Its sign is determined entirely by the overlap geometry of the menus and
their specified winners, independently of all positive rate parameters. Zero covariance occurs
exactly for disjoint menus or for properly nested menus whose larger-menu winner lies outside the
smaller menu. Every pairwise-independent finite collection is then mutually independent: its menus
form a pivot-labelled laminar forest. If the collection uses \(n\) items and has \(r\) maximal
menus, its size is at most \(n-r\); equality is attained exactly by disjoint unions of full flags.
The conclusions extend verbatim to independent proportional-hazards races. The result gives a
parameter-free dependence classification for overlapping choices extracted from one latent
ranking.
\end{abstract}

\section{Introduction}

The Plackett--Luce model assigns a random ranking to alternatives with positive weights. Its
exponential-race representation makes the first-ranked alternative in a menu especially simple:
alternative \(s\) wins menu \(S\) with probability proportional to its weight. Less immediate is
the dependence between two menu choices extracted from the same latent ranking.

This paper determines that dependence exactly for events in which \(s\) wins \(S\). The covariance
sign is structural. Shared competitors create positive covariance, while a nominated winner
appearing as a competitor in the other menu creates negative covariance. The only zero cases are
disjointness and one oriented form of containment. No choice of strictly positive weights creates
an additional cancellation.

The zero-set has a multivariate consequence. Pairwise-independent families form pivot-labelled
laminar forests, and the exponential memoryless property promotes pairwise independence to mutual
independence. A sharp counting theorem then identifies the largest such families.

The model throughout is one random ranking. Choices generated from fresh independent rankings on
different occasions are outside the question considered here.

\subsection*{Relation to earlier work}

Luce's choice model and Plackett's ranking model are classical
\cite{Luce1959,Plackett1975}. The maximal-chain specialization of our independence theorem is the
independence of record indicators in an \(F^\alpha\)-scheme
\cite{BorovkovPfeifer1995}. Decomposability and conditional-independence properties of ranking
models were developed by Critchlow, Fligner, and Verducci, Fishburn, and Csisz\'ar
\cite{CritchlowFlignerVerducci1991,Fishburn1994,Csiszar2009}. Those results explain why nested
constructions can factor. The contribution here is the explicit signed covariance classification,
its exact structural zero-set, the resulting pairwise-to-mutual collapse for arbitrary finite
families, and the sharp laminar-forest bound.

\section{Exponential races and choice events}

Let \(V\) be a finite set. For each \(v\in V\), let
\[
X_v\sim\Exp(\lambda_v),\qquad \lambda_v>0,
\]
independently. For \(A\subseteq V\), write
\[
\lambda(A)=\sum_{v\in A}\lambda_v.
\]

For \(s\in S\subseteq V\), with \(|S|\geq2\), define
\[
C(S,s)=\{X_s<X_v\text{ for every }v\in S\setminus\{s\}\},
\qquad I_{S,s}=\1_{C(S,s)}.
\]
The item \(s\) is the \emph{pivot} of the labelled menu \((S,s)\).

\begin{lemma}[Race lemma]
\label{lem:race}
For nonempty \(A\subseteq V\), the minimum \(M_A=\min_{v\in A}X_v\) has distribution
\(\Exp(\lambda(A))\), and
\[
\Pp(X_a=M_A)=\frac{\lambda_a}{\lambda(A)}.
\]
The identity of the minimizer is independent of \(M_A\). Conditional on \(X_a=t\) and
\(X_v>t\) for all \(v\in A\setminus\{a\}\), the residual variables
\((X_v-t)_{v\in A\setminus\{a\}}\) are independent exponentials with their original rates.
\end{lemma}

\begin{proof}
Independence gives
\[
\Pp(M_A>t)=\prod_{v\in A}e^{-\lambda_vt}=e^{-\lambda(A)t}.
\]
The joint density that \(a\) is the minimizer and \(M_A\in dt\) is
\(\lambda_a e^{-\lambda(A)t}\,dt\). Dividing by the density
\(\lambda(A)e^{-\lambda(A)t}\) proves the winner formula and independence. The final assertion is
the exponential memoryless property applied independently to every surviving clock.
\end{proof}

In particular,
\begin{equation}
\label{eq:marginal}
\Pp(C(S,s))=\frac{\lambda_s}{\lambda(S)}.
\end{equation}

\section{Exact signed covariance}

\begin{theorem}[Signed structural faithfulness]
\label{thm:pair}
Let \((S,a)\) and \((T,b)\) be two labelled menus, each containing at least two items, and put
\(U=S\cup T\), \(\alpha=\lambda_a\), and \(\beta=\lambda_b\). Their covariance is given by the
following exhaustive cases.

\begin{enumerate}
\item If \(a=b\), set
\[
c=\lambda((S\cap T)\setminus\{a\}),\quad
p=\lambda(S\setminus T),\quad
q=\lambda(T\setminus S).
\]
Then
\begin{equation}
\label{eq:same-pivot}
\Cov(I_{S,a},I_{T,a})
=
\frac{\alpha\{c\lambda(U)+pq\}}
     {\lambda(S)\lambda(T)\lambda(U)}
>0.
\end{equation}

\item If \(a\ne b\), \(a\notin T\), and \(b\notin S\), then
\begin{equation}
\label{eq:no-cross}
\Cov(I_{S,a},I_{T,b})
=
\frac{\alpha\beta\lambda(S\cap T)}
     {\lambda(S)\lambda(T)\lambda(U)}
\geq0.
\end{equation}
Equality holds exactly when \(S\cap T=\varnothing\).

\item If \(a\in T\) and \(b\notin S\), then
\begin{equation}
\label{eq:one-cross}
\Cov(I_{S,a},I_{T,b})
=
-\frac{\alpha\beta\lambda(S\setminus T)}
      {\lambda(S)\lambda(T)\lambda(U)}
\leq0.
\end{equation}
Equality holds exactly when \(S\subsetneq T\).

\item If \(b\in S\) and \(a\notin T\), then
\begin{equation}
\label{eq:other-cross}
\Cov(I_{S,a},I_{T,b})
=
-\frac{\alpha\beta\lambda(T\setminus S)}
      {\lambda(S)\lambda(T)\lambda(U)}
\leq0.
\end{equation}
Equality holds exactly when \(T\subsetneq S\).

\item If \(a\in T\) and \(b\in S\), then
\begin{equation}
\label{eq:double-cross}
\Cov(I_{S,a},I_{T,b})
=-\frac{\alpha\beta}{\lambda(S)\lambda(T)}<0.
\end{equation}
\end{enumerate}

Thus the sign, including whether it vanishes, depends only on the two labelled menus and not on the
values of their positive rates.
\end{theorem}

\begin{proof}
If \(a=b\), the two events occur together exactly when \(a\) wins \(U\). Hence
\[
\Pp(C(S,a)\cap C(T,a))=\frac{\alpha}{\lambda(U)}.
\]
Subtracting the product of the marginals in \eqref{eq:marginal}, and using
\[
\lambda(S)\lambda(T)-\alpha\lambda(U)=c\lambda(U)+pq,
\]
gives \eqref{eq:same-pivot}. The right side is strictly positive for non-singleton menus.

Suppose next that \(a\ne b\), \(a\notin T\), and \(b\notin S\). On the joint event, the first clock
in \(U\) must be \(a\) or \(b\). If it is \(a\), the race lemma says that \(b\) subsequently wins
\(T\) with probability \(\beta/\lambda(T)\). The reverse alternative is symmetric. Therefore
\[
\Pp(C(S,a)\cap C(T,b))
=\frac{\alpha}{\lambda(U)}\frac{\beta}{\lambda(T)}
 +\frac{\beta}{\lambda(U)}\frac{\alpha}{\lambda(S)}.
\]
Subtracting the marginal product gives \eqref{eq:no-cross}, because
\[
\lambda(S)+\lambda(T)-\lambda(U)=\lambda(S\cap T).
\]

Suppose \(a\in T\) and \(b\notin S\). The joint event forces \(b\) to be the first clock in \(U\).
After \(b\) rings, the race lemma gives probability \(\alpha/\lambda(S)\) that \(a\) wins \(S\).
Thus
\[
\Pp(C(S,a)\cap C(T,b))
=\frac{\beta}{\lambda(U)}\frac{\alpha}{\lambda(S)}.
\]
Since \(\lambda(U)=\lambda(T)+\lambda(S\setminus T)\), subtraction gives
\eqref{eq:one-cross}. The fourth case is symmetric.

Finally, if \(a\in T\) and \(b\in S\), the first event requires \(X_a<X_b\), while the second
requires \(X_b<X_a\). Their intersection is empty, which proves \eqref{eq:double-cross}.
\end{proof}

\begin{definition}[Pivot-laminar family]
\label{def:pivot-laminar}
A family of labelled menus is pivot-laminar if every pair satisfies exactly one of the following:
their scopes are disjoint; \(S\subsetneq T\) and the pivot of \(T\) lies outside \(S\); or
\(T\subsetneq S\) and the pivot of \(S\) lies outside \(T\).
\end{definition}

\begin{corollary}[No accidental independence]
\label{cor:pair}
Two nontrivial Luce choice events are independent if and only if their labelled menus are
pivot-laminar.
\end{corollary}

\begin{proof}
For indicators, independence is equivalent to zero covariance. The zero cases in
Theorem~\ref{thm:pair} are precisely those in Definition~\ref{def:pivot-laminar}.
\end{proof}

\section{Pairwise independence collapses to mutual independence}

\begin{theorem}[Independence collapse]
\label{thm:collapse}
For every finite family
\[
\mathcal F=\{(S_j,s_j):1\leq j\leq m\}
\]
of nontrivial labelled menus, the following are equivalent:

\begin{enumerate}
\item the events \(C(S_j,s_j)\) are mutually independent;
\item they are pairwise independent;
\item \(\mathcal F\) is pivot-laminar.
\end{enumerate}

When these conditions hold,
\begin{equation}
\label{eq:product}
\Pp\left(\bigcap_{j=1}^m C(S_j,s_j)\right)
=\prod_{j=1}^m\frac{\lambda_{s_j}}{\lambda(S_j)}.
\end{equation}
\end{theorem}

\begin{proof}
Mutual independence implies pairwise independence, and Corollary~\ref{cor:pair} shows that pairwise
independence is equivalent to pivot-laminarity. It remains to prove that pivot-laminarity implies
mutual independence.

The scopes of \(\mathcal F\) form a laminar forest. Distinct maximal scopes are disjoint, so their
subfamilies depend on disjoint collections of clocks. It is enough to treat one maximal scope
\(R\), with pivot \(r\).

Let \(D\) be an intersection of events whose scopes are proper descendants of \(R\). The pivot rule
gives \(r\notin S\) for every such descendant scope \(S\), so \(D\) depends only on relative orders
among clocks in \(R\setminus\{r\}\). For every \(t\geq0\), memorylessness and translation invariance
of comparisons give
\[
\Pp\bigl(D,\ X_v>t\text{ for all }v\in R\setminus\{r\}\bigr)
=\Pp(D)e^{-(\lambda(R)-\lambda_r)t}.
\]
Integrating against the density \(\lambda_r e^{-\lambda_rt}\,dt\) of \(X_r\),
\begin{align*}
\Pp(C(R,r)\cap D)
&=\int_0^\infty
\lambda_r e^{-\lambda_rt}
\Pp\bigl(D,\ X_v>t\text{ for all }v\ne r\bigr)\,dt\\
&=\Pp(D)\frac{\lambda_r}{\lambda(R)}\\
&=\Pp(D)\Pp(C(R,r)).
\end{align*}
Removing the root and applying the same argument recursively down the finite inclusion forest proves
\eqref{eq:product}.
\end{proof}

\section{Sharp extremal structure}

\begin{theorem}[Laminar-forest bound]
\label{thm:bound}
Let \(V_0=\bigcup_{(S,s)\in\mathcal F}S\), let \(n=|V_0|\), and let \(r\) be the number of maximal
scopes in a pairwise-independent family \(\mathcal F\). Then
\[
|\mathcal F|\leq n-r.
\]
Equality holds if and only if every maximal component is a full flag
\[
S_2\subsetneq S_3\subsetneq\cdots\subsetneq S_k,
\qquad |S_i|=i,
\]
where the pivot of \(S_i\) is its unique element outside \(S_{i-1}\) for \(i\geq3\), and the pivot
of \(S_2\) is one of its two elements.
\end{theorem}

\begin{proof}
All pivots are distinct. Indeed, disjoint scopes have distinct pivots, while for nested scopes the
pivot of the larger scope lies outside the smaller one.

For each maximal component, choose a minimal scope \(L\) and an element \(x\in L\) other than its
pivot. This \(x\) cannot be the pivot of another scope. Any scope containing \(x\) intersects \(L\),
so laminarity makes the two scopes comparable. A proper subscope contradicts the minimality of
\(L\), equality of scopes is excluded by pivot-laminarity, and a proper superscope must have its
pivot outside \(L\). Thus every maximal component
contains a distinct nonpivot. At most \(n-r\) elements can therefore serve as pivots, proving the
bound.

If equality holds, each component has exactly one nonpivot. Every leaf of its inclusion tree
contains a nonpivot by the preceding argument, so the tree has one leaf and is a chain. Its smallest
scope has at least two elements, and every strict inclusion adds at least one element. Equality
forces the smallest scope to have size two and every later inclusion to add exactly the new pivot.
This is a full flag. Conversely, a disjoint union of such flags has one nonpivot per component and
attains the bound.
\end{proof}

\begin{corollary}[One-component extremizers]
On \(n\) labelled items, a pairwise-independent family has at most \(n-1\) events. The families of
size \(n-1\) are in bijection with the \(n!\) permutations of the items.
\end{corollary}

\begin{proof}
For a full flag, list first its unique nonpivot and then, in increasing scope order, the pivot added
at each stage. This gives a permutation. Conversely, a permutation
\((v_1,\ldots,v_n)\) defines \(S_k=\{v_1,\ldots,v_k\}\) with pivot \(v_k\) for
\(2\leq k\leq n\).
\end{proof}

\section{Proportional-hazards extension}

\begin{corollary}
Suppose the independent clocks \(T_v\) have survival functions
\[
\Pp(T_v>t)=\exp[-\lambda_v H(t)],
\]
where \(H:[0,\infty)\to[0,\infty)\) is continuous and strictly increasing, \(H(0)=0\), and
\(H(t)\to\infty\). Every conclusion of Theorems~\ref{thm:pair}, \ref{thm:collapse}, and
\ref{thm:bound} remains valid.
\end{corollary}

\begin{proof}
The transformed variables \(H(T_v)\) are independent \(\Exp(\lambda_v)\) variables. Since \(H\) is
strictly increasing, it preserves every choice event.
\end{proof}

\section{Scope and implications}

Theorem~\ref{thm:pair} supplies the exact pairwise dependence graph of top-choice indicators derived
from one Plackett--Luce ranking. In particular, its sign pattern is invariant over the full positive
parameter space. When rank data are broken into overlapping menu-choice features, the menu geometry
alone identifies which features factor exactly and whether the remaining pairwise dependence is
positive or negative.

The conclusions have precise limits. They concern specified-winner indicators, not arbitrary
functions of a ranking or arbitrary conditional-independence statements. Strictly positive rates
and non-singleton menus exclude deterministic events. The proportional-hazards corollary uses a
common baseline; unrelated clock distributions need not satisfy the same nested independence.

\begin{thebibliography}{9}

\bibitem{BorovkovPfeifer1995}
K.~A. Borovkov and D. Pfeifer.
\newblock On record indices and record times.
\newblock \emph{Journal of Statistical Planning and Inference}, 1995.
\newblock \url{https://doi.org/10.1016/0378-3758(94)00063-8}.

\bibitem{CritchlowFlignerVerducci1991}
D.~E. Critchlow, M.~A. Fligner, and J.~S. Verducci.
\newblock Probability models on rankings.
\newblock \emph{Journal of Mathematical Psychology}, 1991.

\bibitem{Csiszar2009}
V. Csisz\'ar.
\newblock On L-decomposability of random orderings.
\newblock \emph{Journal of Mathematical Psychology}, 2009.
\newblock \url{https://doi.org/10.1016/j.jmp.2009.04.011}.

\bibitem{Fishburn1994}
P.~C. Fishburn.
\newblock On choice probabilities derived from ranking distributions.
\newblock \emph{Journal of Mathematical Psychology}, 1994.
\newblock \url{https://doi.org/10.1006/jmps.1994.1017}.

\bibitem{Luce1959}
R.~D. Luce.
\newblock \emph{Individual Choice Behavior: A Theoretical Analysis}.
\newblock Wiley, 1959.

\bibitem{Plackett1975}
R.~L. Plackett.
\newblock The analysis of permutations.
\newblock \emph{Applied Statistics}, 1975.
\newblock \url{https://doi.org/10.2307/2346567}.

\end{thebibliography}

\end{document}
